% GENERATED FILE. Edit canonical lessons, then run npm run generate:books.
% canonical-source-hash: fnv1a64:2e44c791cb6afb3e
% canonical-lessons: TE-C356-utsavam, TE-C356-sanmanincu, TE-C356-grhapravesam, TE-C356-vidkolu, TE-C356-dattata

\chapter{Celebrations}
\label{ch:te-celebrations}
\hlchaptermodality{356}

\begin{chapteropening}
\textbf{By the end of this chapter:} \emph{I can say a public celebration, to felicitate, a housewarming, a farewell and adoption.}

Complete the last lesson of chapter 356: I can say a public celebration, to felicitate, a housewarming, a farewell and adoption.
\end{chapteropening}

\section[\texorpdfstring{\te{ఉత్సవం} (utsavaṁ)}{ఉత్సవం (utsavaṁ)}]{\te{ఉత్సవం} (utsavaṁ) — a public celebration; a festival}

\label{lesson:TE-C356-utsavam}



\begin{quote}
Before the new one: say the Telugu for a birth anniversary, then the Telugu for a felicitation.
\end{quote}

\subsection*{You'll want to know: \te{ఉత్సవం}}
\textbf{\te{ఉత్సవం}} — \emph{utsavaṁ} — \textquotedblleft{}a public celebration; a festival\textquotedblright{}.

\begin{quote}
\te{పట్టణంలో} \te{పెద్ద} \te{ఉత్సవం} \te{జరిగింది}.
\end{quote}

\begin{quote}
\emph{paṭṭaṇaṁlō pedda utsavaṁ jarigindi.}
\end{quote}

\begin{quote}
A big celebration took place in the town.
\end{quote}

Plural: \te{ఉత్సవాలు}, \emph{utsavālu}.

\begin{cousinweb}
\textbf{\te{ఉత్సవం}} (\emph{utsavaṁ}) is borrowed from Sanskrit \emph{utsava}, \textquotedblleft{}an enterprise\textquotedblright{}, and later \textquotedblleft{}a festival\textquotedblright{} (CDIAL 1876).
\end{cousinweb}

\begin{grammarlens}[title={What you've built}]
One more word for telling what happened.
\end{grammarlens}

\subsection*{Your turn}
\begin{itemize}
\raggedright
  \item \emph{Say it:} \emph{utsavaṁ}
  \item \emph{Say it:} \emph{utsavaṁ}, once more
  \item \emph{Say it:} \emph{paṭṭaṇaṁlō pedda utsavaṁ jarigindi.} — A big celebration took place in the town.
  \item \emph{Recall:} say the Telugu for a birth anniversary, then the Telugu for a felicitation, then say \emph{utsavaṁ} again
\end{itemize}

\begin{tcolorbox}[breakable,colback=teal!4,colframe=teal!35!black,title={Before you move on}]
What does \te{ఉత్సవం} mean? (\textquotedblleft{}A public celebration; a festival\textquotedblright{}.) Say it once more.
\end{tcolorbox}
\section[\texorpdfstring{\te{సన్మానించు} (sanmāniñcu)}{సన్మానించు (sanmāniñcu)}]{\te{సన్మానించు} (sanmāniñcu) — to felicitate; to honour}

\label{lesson:TE-C356-sanmanincu}



\begin{quote}
Before the new one: say the Telugu for a felicitation, then the Telugu for a public celebration.
\end{quote}

\subsection*{You'll want to know: \te{సన్మానించు}}
\textbf{\te{సన్మానించు}} — \emph{sanmāniñcu} — \textquotedblleft{}to felicitate; to honour\textquotedblright{}.

\begin{quote}
\te{పాఠశాల} \te{తన} \te{పాత} \te{విద్యార్థులను} \te{సన్మానించింది}.
\end{quote}

\begin{quote}
\emph{pāṭhaśāla tana pāta vidyārthulanu sanmāniñcindi.}
\end{quote}

\begin{quote}
The school honoured its former students.
\end{quote}

Past (I …): \te{సన్మానించాను}, \emph{sanmāniñcānu}.

The verb of \te{సన్మానం}.

\begin{cousinweb}
\textbf{\te{సన్మానించు}} (\emph{sanmāniñcu}) is built on Sanskrit \emph{sanmāna}, \textquotedblleft{}respect, esteem\textquotedblright{} (Monier-Williams), with the Telugu verb ending \emph{-iñcu}.
\end{cousinweb}

\begin{grammarlens}[title={What you've built}]
One more word for telling what happened.
\end{grammarlens}

\subsection*{Your turn}
\begin{itemize}
\raggedright
  \item \emph{Say it:} \emph{sanmāniñcu}
  \item \emph{Say it:} \emph{sanmāniñcu}, once more
  \item \emph{Say it:} \emph{pāṭhaśāla tana pāta vidyārthulanu sanmāniñcindi.} — The school honoured its former students.
  \item \emph{Recall:} say the Telugu for a felicitation, then the Telugu for a public celebration, then say \emph{sanmāniñcu} again
\end{itemize}

\begin{tcolorbox}[breakable,colback=teal!4,colframe=teal!35!black,title={Before you move on}]
What does \te{సన్మానించు} mean? (\textquotedblleft{}To felicitate; to honour\textquotedblright{}.) Say it once more.
\end{tcolorbox}
\section[\texorpdfstring{\te{గృహప్రవేశం} (gṛhapravēśaṁ)}{గృహప్రవేశం (gṛhapravēśaṁ)}]{\te{గృహప్రవేశం} (gṛhapravēśaṁ) — a housewarming}

\label{lesson:TE-C356-grhapravesam}



\begin{quote}
Before the new one: say the Telugu for a public celebration, then the Telugu for to felicitate.
\end{quote}

\subsection*{You'll want to know: \te{గృహప్రవేశం}}
\textbf{\te{గృహప్రవేశం}} — \emph{gṛhapravēśaṁ} — \textquotedblleft{}a housewarming\textquotedblright{}.

\begin{quote}
\te{మా} \te{కొత్త} \te{ఇంటి} \te{గృహప్రవేశం} \te{ఆదివారం} \te{జరిగింది}.
\end{quote}

\begin{quote}
\emph{mā kotta iṇṭi gṛhapravēśaṁ ādivāraṁ jarigindi.}
\end{quote}

\begin{quote}
The housewarming of our new house was on Sunday.
\end{quote}

\te{గృహం}, a house, and a Sanskrit word for entering: the first entry into a new home.

\begin{cousinweb}
\textbf{\te{గృహప్రవేశం}} (\emph{gṛhapravēśaṁ}) is borrowed from Sanskrit: \emph{gṛha}, \textquotedblleft{}a house\textquotedblright{}, and \emph{praveśa}, \textquotedblleft{}an entering\textquotedblright{} (Monier-Williams).
\end{cousinweb}

\begin{grammarlens}[title={What you've built}]
One more word for telling what happened.
\end{grammarlens}

\subsection*{Your turn}
\begin{itemize}
\raggedright
  \item \emph{Say it:} \emph{gṛhapravēśaṁ}
  \item \emph{Say it:} \emph{gṛhapravēśaṁ}, once more
  \item \emph{Say it:} \emph{mā kotta iṇṭi gṛhapravēśaṁ ādivāraṁ jarigindi.} — The housewarming of our new house was on Sunday.
  \item \emph{Recall:} say the Telugu for a public celebration, then the Telugu for to felicitate, then say \emph{gṛhapravēśaṁ} again
\end{itemize}

\begin{tcolorbox}[breakable,colback=teal!4,colframe=teal!35!black,title={Before you move on}]
What does \te{గృహప్రవేశం} mean? (\textquotedblleft{}A housewarming\textquotedblright{}.) Say it once more.
\end{tcolorbox}
\section[\texorpdfstring{\te{వీడ్కోలు} (vīḍkōlu)}{వీడ్కోలు (vīḍkōlu)}]{\te{వీడ్కోలు} (vīḍkōlu) — a farewell}

\label{lesson:TE-C356-vidkolu}



\begin{quote}
Before the new one: say the Telugu for to felicitate, then the Telugu for a housewarming.
\end{quote}

\subsection*{You'll want to know: \te{వీడ్కోలు}}
\textbf{\te{వీడ్కోలు}} — \emph{vīḍkōlu} — \textquotedblleft{}a farewell\textquotedblright{}.

\begin{quote}
\te{ఆఫీసులో} \te{అతనికి} \te{ఘనంగా} \te{వీడ్కోలు} \te{ఇచ్చారు}.
\end{quote}

\begin{quote}
\emph{āphīsulō ataniki ghanaṅgā vīḍkōlu iccāru.}
\end{quote}

\begin{quote}
At the office they gave him a grand farewell.
\end{quote}

\te{వీడ్కోలు} \te{ఇచ్చు}, to see someone off.

\begin{cousinweb}
\textbf{\te{వీడ్కోలు}} (\emph{vīḍkōlu}) is a Telugu word for a send-off; its origin is not settled.
\end{cousinweb}

\begin{grammarlens}[title={What you've built}]
One more word for telling what happened.
\end{grammarlens}

\subsection*{Your turn}
\begin{itemize}
\raggedright
  \item \emph{Say it:} \emph{vīḍkōlu}
  \item \emph{Say it:} \emph{vīḍkōlu}, once more
  \item \emph{Say it:} \emph{āphīsulō ataniki ghanaṅgā vīḍkōlu iccāru.} — At the office they gave him a grand farewell.
  \item \emph{Recall:} say the Telugu for to felicitate, then the Telugu for a housewarming, then say \emph{vīḍkōlu} again
\end{itemize}

\begin{tcolorbox}[breakable,colback=teal!4,colframe=teal!35!black,title={Before you move on}]
What does \te{వీడ్కోలు} mean? (\textquotedblleft{}A farewell\textquotedblright{}.) Say it once more.
\end{tcolorbox}
\section[\texorpdfstring{\te{దత్తత} (dattata)}{దత్తత (dattata)}]{\te{దత్తత} (dattata) — adoption}

\label{lesson:TE-C356-dattata}



\begin{quote}
Before the new one: say the Telugu for a housewarming, then the Telugu for a farewell.
\end{quote}

\subsection*{You'll want to know: \te{దత్తత}}
\textbf{\te{దత్తత}} — \emph{dattata} — \textquotedblleft{}adoption\textquotedblright{}.

\begin{quote}
\te{వాళ్ళు} \te{ఒక} \te{పాపను} \te{దత్తత} \te{తీసుకున్నారు}.
\end{quote}

\begin{quote}
\emph{vāḷḷu oka pāpanu dattata tīsukunnāru.}
\end{quote}

\begin{quote}
They adopted a baby girl.
\end{quote}

\te{దత్తత} \te{తీసుకో}, to adopt.

\begin{cousinweb}
\textbf{\te{దత్తత}} (\emph{dattata}) is built on Sanskrit \emph{datta}, \textquotedblleft{}given\textquotedblright{} (CDIAL 6140): a child given into a new family.
\end{cousinweb}

\begin{grammarlens}[title={What you've built}]
One more word for telling what happened.
\end{grammarlens}

\subsection*{Your turn}
\begin{itemize}
\raggedright
  \item \emph{Say it:} \emph{dattata}
  \item \emph{Say it:} \emph{dattata}, once more
  \item \emph{Say it:} \emph{vāḷḷu oka pāpanu dattata tīsukunnāru.} — They adopted a baby girl.
  \item \emph{Recall:} say the Telugu for a housewarming, then the Telugu for a farewell, then say \emph{dattata} again
\end{itemize}

\begin{tcolorbox}[breakable,colback=teal!4,colframe=teal!35!black,title={Before you move on}]
What does \te{దత్తత} mean? (\textquotedblleft{}Adoption\textquotedblright{}.) Say it once more.
\end{tcolorbox}
